\subsubsection{Velocidades linear e angular}
\label{sec:dinAcoplada_veloc_linear_angular}

As velocidades linear ${}^{B}\!\vec v_{G_i}$ e angular ${}^{B}\!\vec\omega_i$
do elo $i$ são obtidas a partir do jacobiano geométrico avaliado no respectivo
centro de massa \cite{craig2005,Wang2003KinematicsDynamics}.
Escrevendo o jacobiano na forma espacial \cite{murray1994mathematical},
tem-se um operador $J_{\text{sp},i}(\vec q\, \, )$ que
relaciona diretamente as velocidades generalizadas
$\dot{\vec q\, \, }$ com as velocidades cinemáticas empregadas na energia
cinética:
\begin{equation}
\begin{bmatrix}
{}^{B}\!\vec v_{G_i} \\
{}^{B}\!\vec\omega_i
\end{bmatrix}
=
J_{\text{sp},i}(\vec q\,)\,\dot{\vec q\,}.
\label{eq:jacobianos_vGi1}
\end{equation}

Em seguida, decompõe-se o jacobiano espacial em parte linear e parte angular

\begin{equation}
J_{\text{sp},i}(\vec q\,)
=
\begin{bmatrix}
J_{v,i}(\vec q\,) \\
J_{\omega,i}(\vec q\,)
\end{bmatrix},
\label{eq:jacobianos_vGi2}
\end{equation}

onde o termo $J_{v,i}(\vec q\, )$ é a parte linear, enquanto
$J_{\omega,i}(\vec q\, )$ é a parte angular do jacobiano espacial do elo $i$.

Como discutido em \cite{sciavicco_siciliano2010}, no caso de juntas
revolutas as colunas de $J_{v,i}$ e $J_{\omega,i}$ associadas às coordenadas
articulares podem ser escritas a partir dos eixos de junta e das posições
dos centros de massa.
Se ${}^{B}\!\vec z_j$ denota o eixo de rotação da junta $j$ em $\{B\}$,
${}^{B}\!\vec r_{O_j}$ a posição da sua origem e ${}^{B}\!\vec r_{G_i}$ a
posição do centro de massa do elo $i$, então, para uma junta revoluta $j$
que atua sobre o elo $i$, a coluna correspondente é

\begin{equation}
    J_{v,i}^{(j)}(\vec q\, )
    =
    {}^{B}\!\vec z_j \times
    \bigl({}^{B}\!\vec r_{G_i}(\vec q\, ) - {}^{B}\!\vec r_{O_j}(\vec q\, )\bigr),
    \qquad
    J_{\omega,i}^{(j)}(\vec q\, )
    =
    {}^{B}\!\vec z_j,
\end{equation}
para $j\le i$, ao passo que, para $j>i$, essas colunas são nulas,
\begin{equation}
J_{v,i}^{(j)}(\vec q\, ) = \vec 0,
\qquad
J_{\omega,i}^{(j)}(\vec q\, ) = \vec 0.
\end{equation}
Essas expressões coincidem com as de um manipulador de base fixa, pois se
referem apenas às coordenadas articulares $\theta_j$; as colunas de
$J_{\text{sp},i}(\vec q\, )$ associadas às coordenadas de atitude da base são
determinadas pela cinemática da base e aparecem de forma explícita nas
relações entre $\dot{\vec\eta}_B$, ${}^{B}\!\vec\omega_B$ e na construção de
$M(\vec q\, )$.

%%%%%%%%%%==========%%%%%%%%%%==========%%%%%%%%%%==========%%%%%%%%%%==========%%%%%%%%%%==========

A expressão de Coriolis e centrífuga para cada elo decorre diretamente da
energia cinética escrita no espaço de velocidades espaciais.
Como em \cite{featherstone2008}, para o elo $i$ a energia cinética pode ser
escrita como
\begin{equation}
T_i
=
\frac{1}{2}
\left(
J_{\text{sp},i}(\vec q\, )\,\dot{\vec q\, }
\right)^{T}
M_{i,\text{spatial}}
\left(
J_{\text{sp},i}(\vec q\, )\,\dot{\vec q\, }
\right),
\label{eq:Ti_spatial_energy}
\end{equation}
onde $M_{i,\text{spatial}}$ é a matriz espacial de
inércia do elo, composta pela massa e pelo tensor de inércia do elo escritos no
formalismo de velocidades espaciais 6D.
Ademais, a derivada de \eqref{eq:Ti_spatial_energy} em relação às velocidades generalizadas produz
\begin{equation}
\frac{\partial T_i}{\partial \dot{\vec q\, }}
=
J_{\text{sp},i}^{T}\,
M_{i,\text{spatial}}\,
J_{\text{sp},i}\,\dot{\vec q\, },
\end{equation}
e, por sua vez, a derivada temporal gera o termo de Coriolis e centrífugo:
\begin{equation}
\frac{d}{dt}
\left(
\frac{\partial T_i}{\partial \dot{\vec q\, }}
\right)
=
J_{\text{sp},i}^{T}
M_{i,\text{spatial}}
\left(
\dot J_{\text{sp},i}\,\dot{\vec q\, }
\right)
+
\left(\dot J_{\text{sp},i}\right)^{T}
M_{i,\text{spatial}}
J_{\text{sp},i}\,\dot{\vec q\, }.
\label{eq:derivada_spatial}
\end{equation}

Como $M_{i,\text{spatial}}$ é constante para cada elo (sua massa e seu tensor de
inércia são definidos no referencial do elo e independentemente de $\vec q\, $),
o termo simétrico
$(\dot J_{\text{sp},i})^{T}M_{i,\text{spatial}}J_{\text{sp},i}\dot{\vec q\, }$
pode ser incorporado à parte inercial da dinâmica.
Assim, a parcela do elo
$i$ que efetivamente compõe o vetor de forças generalizadas de Coriolis é
\begin{equation}
C_i(\vec q\, ,\dot{\vec q\, })\,\dot{\vec q\, }
=
J_{\text{sp},i}^{T}\,
M_{i,\text{spatial}}
\left(
\dot J_{\text{sp},i}\,\dot{\vec q\, }
\right).
\label{eq:termos_centrif_coriolis}
\end{equation}

A expressão \eqref{eq:termos_centrif_coriolis} é equivalente à forma
clássica dos termos de Coriolis e centrífugos escrita a partir dos símbolos de
Christoffel \cite{featherstone2008}.
Para um manipulador com matriz de inércia
$B(\vec q\, )$,
cujos elementos são $b_{ij}(\vec q\, )$,
define-se o tensor de coeficientes, conforme \cite{sciavicco_siciliano2010} apresenta:
\begin{equation}
c_{ijk}(\vec q\, )
=
\frac{1}{2}
\left(
\frac{\partial b_{ij}}{\partial q_k}
+
\frac{\partial b_{ik}}{\partial q_j}
-
\frac{\partial b_{jk}}{\partial q_i}
\right),
\label{eq:christoffel_coefs}
\end{equation}
conhecidos como símbolos de Christoffel de primeira espécie \cite{sciavicco_siciliano2010}.
A matriz de Coriolis é então construída como
\begin{equation}
c_{ij}(\vec q\, ,\dot{\vec q\, })
=
\sum_{k=1}^{n} c_{ijk}(\vec q\, )\,\dot q_k,
\qquad
i,j = 1,\dots,n,
\label{eq:C_Christoffel}
\end{equation}
de modo que o vetor de forças generalizadas associado a Coriolis e efeitos
centrífugos assume a forma
\begin{equation}
\bigl[C(\vec q\, ,\dot{\vec q\, })\,\dot{\vec q\, }\bigr]_i
=
\sum_{j=1}^{n} c_{ij}(\vec q\, ,\dot{\vec q\, })\,\dot q_j
=
\sum_{j=1}^{n}\sum_{k=1}^{n}
c_{ijk}(\vec q\, )\,\dot q_j \dot q_k.
\label{eq:C_Christoffel_vetor}
\end{equation}

Tanto a representação
\eqref{eq:christoffel_coefs}--\eqref{eq:C_Christoffel_vetor},
quanto a forma espacial em \eqref{eq:termos_centrif_coriolis} derivam da mesma
estrutura da energia cinética.
A diferença está na escolha das variáveis intermediárias:
na formulação clássica, parte-se diretamente da
matriz de inércia $B(\vec q\, )$ no espaço das coordenadas generalizadas;
na formulação espacial, a energia cinética de cada elo é escrita
como \eqref{eq:Ti_spatial_energy} e toda a dependência em
$(\vec q\, ,\dot{\vec q\, })$ é encapsulada nos jacobianos espaciais $J_{\text{sp},i}(\vec q\, )$.

%%%%%%%%%%==========%%%%%%%%%%==========%%%%%%%%%%==========%%%%%%%%%%==========%%%%%%%%%%==========

Em ambiente de microgravidade, o centro de massa do sistema praticamente
coincide com o centro de gravidade e o torque gravitacional interno em torno
do centro de massa é desprezível \cite{wie2008space}.
Além disso, para uma órbita circular, a contribuição da energia potencial
gravitacional na dinâmica de atitude é praticamente constante.
Assim, no modelo adotado, toma-se a energia potencial generalizada como
$V(\vec q\, ) = 0$ e a Lagrangiana reduz-se
a $L(\vec q\, ,\dot{\vec q\, }) = T(\vec q\, ,\dot{\vec q\, })$.
O conjunto espaçonave--\gls{mr} é, então, descrito diretamente pelas equações
de Lagrange em sua forma padrão \cite{murray1994mathematical}:
\begin{equation}
M(\vec q\, )\,\ddot{\vec q\, }
+
C(\vec q\, ,\dot{\vec q\, })\,\dot{\vec q\, }
=
\vec\tau,
\label{eq:lagrange_formato_padrao}
\end{equation}
em que $\vec q\, $ é o vetor de coordenadas generalizadas, $M(\vec q\, )$ é a
matriz de inércia generalizada e $C(\vec q\, ,\dot{\vec q\, })\,\dot{\vec q\, }$
agrega os termos de Coriolis e centrífugos.
O foco desta seção é a construção de $M(\vec q\, )$ a partir das
contribuições de massa e inércia da base e dos elos do manipulador.

As coordenadas generalizadas são organizadas de modo a reunir, em um único
vetor, os parâmetros de atitude da base e as coordenadas articulares do
manipulador:
\begin{equation}
\vec q\, 
=
\begin{bmatrix}
\vec\eta_B \\
\vec\theta
\end{bmatrix},
\qquad
\dot{\vec q\, }
=
\begin{bmatrix}
\dot{\vec\eta}_B \\
\dot{\vec\theta}
\end{bmatrix},
\label{eq:q_generalizado_base_mr}
\end{equation}
em que $\vec\eta_B $ representa os ângulos de Euler da espaçonave e
$\vec\theta$ contém as coordenadas articulares do manipulador.
A relação cinemática entre as taxas dos ângulos de Euler e a velocidade
angular do corpo é dada pelas equações cinemáticas diferenciais
\cite{Schaub2010AttitudeDynamics}, que, para a sequência 3--2--1,
podem ser escritas no formato
$\dot{\vec\eta}_B = B(\vec\eta_B)\,{}^{B}\!\vec\omega_B$.
Neste trabalho, utiliza-se a forma inversa,
\begin{equation}
    {}^{B}\!\vec\omega_B
    =
    \mathbf{T}_\omega(\vec\eta_B)\,\dot{\vec\eta}_B,
    \qquad
    \mathbf{T}_\omega(\vec\eta_B)=B^{-1}(\vec\eta_B),
\end{equation}
de modo a expressar a velocidade angular da base em função das
velocidades generalizadas de atitude.
Ao reunir a velocidade angular da espaçonave e as velocidades articulares em um
único vetor cinemático, obtém-se
\begin{equation}
    \dot{\vec y}
    =
    \begin{bmatrix}
    {}^{B}\!\vec\omega_B \\
    \dot{\vec\theta}
    \end{bmatrix}
    =
    \underbrace{
    \begin{bmatrix}
    \mathbf{T}_\omega(\vec\eta_B) & 0 \\
    0 & I_{n}
    \end{bmatrix}}_{\displaystyle \mathbf{J}_y(\vec q\, )}
    \dot{\vec q\, },
\end{equation}
isto é, as velocidades físicas relevantes são combinações lineares das
velocidades generalizadas $\dot{\vec q\, }$.
Com essa definição, a energia cinética total do conjunto pode ser escrita
na forma matricial
\begin{equation}
T =
\frac{1}{2}\,\dot{\vec q\, }^{\,T}\,M(\vec q\, )\,\dot{\vec q\, },
\label{eq:T_matriz_corrigida}
\end{equation}
onde $M(\vec q\, )$ é a matriz de inércia generalizada do sistema
espaçonave--\gls{mr}. Essa matriz é construída somando-se as contribuições
de cada elo, obtidas a partir de seus jacobianos espaciais, e a
contribuição da base física, conforme detalhado a seguir.

Para cada elo, a matriz $M_{i,\text{spatial}}$ é um operador de inércia
espacial $6\times 6$ que combina, em um único bloco, os efeitos
translacionais e rotacionais em torno do centro de massa expresso em
$\{B\}$.
Se ${}^{B}\!p_{G_i}$ denota a posição do centro de massa do elo em relação à
origem de $\{B\}$ e ${}^{B}\!I_i$ o tensor de inércia do elo em torno desse
centro, a forma padrão do operador de inércia espacial é

\begin{equation}
M_{i,\text{spatial}}
=
\begin{bmatrix}
m_i I_3 & -m_i\,S({}^{B}\!p_{G_i})\\[2pt]
m_i\,S({}^{B}\!p_{G_i}) &
{}^{B}\!I_i - m_i\,S({}^{B}\!p_{G_i})\,S({}^{B}\!p_{G_i})^{T}
\end{bmatrix},
\end{equation}

em que $m_i$ é a massa do $i$-ésimo elo, $I_3$ é a matriz identidade $3\times 3$,
${}^{B}\!p_{G_i}$ é o vetor posição do centro de massa do elo $i$
em relação à origem do referencial $\{B\}$, ${}^{B}\!I_i$
é o tensor de inércia do elo $i$ em torno do seu centro de massa, expresso em
$\{B\}$, e $S(\cdot)$ é o operador matricial antissimétrico associado ao produto
vetorial, tal que $S(\vec a)\,\vec b = \vec a \times \vec b$ para quaisquer
$\vec a,\vec b$.

De maneira análoga, a espaçonave (base) é descrita por um operador de inércia
espacial $M_{\text{base},\text{spatial}}$.
No formalismo de velocidades espaciais \cite{featherstone2008}, uma matriz de
inércia espacial relaciona a velocidade espacial
$\vec v_{\text{sp}} = [\,{}^{B}\!\vec v^{\,T}\;{}^{B}\!\vec\omega^{\,T}\,]^{T}$
com a quantidade de movimento espacial
$\vec h_{\text{sp}}$ por meio de
$\vec h_{\text{sp}} = M_{\text{spatial}}\,\vec v_{\text{sp}}$.

Para a base, $M_{\text{base},\text{spatial}}$ é construído a partir da massa
$m_B$ e do tensor ${}^{B}\!I_B$ em torno do seu centro de massa, em completa
analogia com $M_{i,\text{spatial}}$ definido para os elos.

A energia cinética total do conjunto espaçonave--\gls{mr} resulta da soma das
energias cinéticas de cada elo e da base.
Utilizando a forma espacial em \eqref{eq:Ti_spatial_energy} para cada elo,
tem-se

\begin{equation}
T
=
\sum_{i=1}^{n}
\frac{1}{2}
\bigl(J_{\text{sp},i}(\vec q\, )\,\dot{\vec q\, }\bigr)^{T}
M_{i,\text{spatial}}
\bigl(J_{\text{sp},i}(\vec q\, )\,\dot{\vec q\, }\bigr)
+
\frac{1}{2}\,
\vec v_{\text{base,sp}}^{\,T}\,
M_{\text{base},\text{spatial}}\,
\vec v_{\text{base,sp}},
\end{equation}

em que $\vec v_{\text{base,sp}}$ é a velocidade espacial da base, escrita
também como combinação linear das velocidades generalizadas
$\dot{\vec q\, }$.
Comparando essa expressão com a forma quadrática de \eqref{eq:T_matriz_corrigida}, obtém-se

\begin{equation}
T =
\frac{1}{2}\,\dot{\vec q\, }^{\,T}\,M(\vec q\, )\,\dot{\vec q\, },
\end{equation}

que é a matriz de inércia generalizada $M(\vec q\, )$ é dada pela soma das
contribuições espaciais de todos os elos, ponderadas pelos respectivos
jacobianos espaciais, acrescida da contribuição inercial da base:

\begin{equation}
M(\vec q\, ) =
\sum_{i=1}^{n}
J_{\text{sp},i}^{T}(\vec q\, )\,
M_{i,\text{spatial}}\,
J_{\text{sp},i}(\vec q\, )
+
M_{\text{base},\text{spatial}}.
\label{eq:M_total_formula_corrigida}
\end{equation}

Cada termo
$J_{\text{sp},i}^{T}(\vec q\, \, \, )\,M_{i,\text{spatial}}\,J_{\text{sp},i}(\vec q\, \, \, )$
é, portanto, a contribuição espacial do elo $i$ para a matriz de inércia
generalizada.
Nessa expressão, cada jacobiano espacial $J_{\text{sp},i}(\vec q\, \, )$ mapeia as
velocidades generalizadas $\dot{\vec q\, \, }$ na velocidade espacial (linear e
angular) do centro de massa do elo $i$ em $\{B\}$, garantindo consistência
entre a cinemática diferencial e a estrutura de $M(\vec q\, \, )$.
